\documentclass[10pt,a4paper]{amsart}
\usepackage{amsmath,amssymb,amsthm,graphicx}
\usepackage{verbatim}
\newenvironment{reference}{\comment}{\endcomment}
%\newenvironment{reference}{}{}
\newenvironment{slogan}{\comment}{\endcomment}
\newenvironment{history}{\comment}{\endcomment}

% For commutative diagrams we use Xy-pic
\usepackage[all]{xy}

% We use 2cell for 2-commutative diagrams.
\xyoption{2cell}
% Theorem environments.
%
\theoremstyle{plain}
\newtheorem{theorem}[subsection]{Theorem}
\newtheorem{proposition}[subsection]{Proposition}
\newtheorem{lemma}[subsection]{Lemma}

\theoremstyle{definition}
\newtheorem{definition}[subsection]{Definition}
\newtheorem{example}[subsection]{Example}
\newtheorem{exercise}[subsection]{Exercise}
\newtheorem{situation}[subsection]{Situation}

\theoremstyle{remark}
\newtheorem{remark}[subsection]{Remark}
\newtheorem{remarks}[subsection]{Remarks}

\numberwithin{equation}{subsection}

% Macros
%
\def\lim{\mathop{\mathrm{lim}}\nolimits}
\def\colim{\mathop{\mathrm{colim}}\nolimits}
\def\Spec{\mathop{\mathrm{Spec}}}
\def\Hom{\mathop{\mathrm{Hom}}\nolimits}
\def\Ext{\mathop{\mathrm{Ext}}\nolimits}
\def\SheafHom{\mathop{\mathcal{H}\!\mathit{om}}\nolimits}
\def\SheafExt{\mathop{\mathcal{E}\!\mathit{xt}}\nolimits}
\def\Sch{\mathit{Sch}}
\def\Mor{\mathop{\mathrm{Mor}}\nolimits}
\def\Ob{\mathop{\mathrm{Ob}}\nolimits}
\def\Sh{\mathop{\mathit{Sh}}\nolimits}
\def\NL{\mathop{N\!L}\nolimits}
\def\CH{\mathop{\mathrm{CH}}\nolimits}
\def\proetale{{pro\text{-}\acute{e}tale}}
\def\etale{{\acute{e}tale}}
\def\QCoh{\mathit{QCoh}}
\def\Ker{\mathop{\mathrm{Ker}}}
\def\Im{\mathop{\mathrm{Im}}}
\def\Coker{\mathop{\mathrm{Coker}}}
\def\Coim{\mathop{\mathrm{Coim}}}

% Boxtimes
%
\DeclareMathSymbol{\boxtimes}{\mathbin}{AMSa}{"02}

%
% Macros for moduli stacks/spaces
%
\def\QCohstack{\mathcal{QC}\!\mathit{oh}}
\def\Cohstack{\mathcal{C}\!\mathit{oh}}
\def\Spacesstack{\mathcal{S}\!\mathit{paces}}
\def\Quotfunctor{\mathrm{Quot}}
\def\Hilbfunctor{\mathrm{Hilb}}
\def\Curvesstack{\mathcal{C}\!\mathit{urves}}
\def\Polarizedstack{\mathcal{P}\!\mathit{olarized}}
\def\Complexesstack{\mathcal{C}\!\mathit{omplexes}}
% \Pic is the operator that assigns to X its picard group, usage \Pic(X)
% \Picardstack_{X/B} denotes the Picard stack of X over B
% \Picardfunctor_{X/B} denotes the Picard functor of X over B
\def\Pic{\mathop{\mathrm{Pic}}\nolimits}
\def\Picardstack{\mathcal{P}\!\mathit{ic}}
\def\Picardfunctor{\mathrm{Pic}}
\def\Deformationcategory{\mathcal{D}\!\mathit{ef}}
\title{Two choices of splitting give homotopic connecting morphisms}
\author{The Stacks Project Authors}
\date{stacks-project:homology.tex:local-snapshot-20261012; GNU FDL 1.2 or later}
\begin{document}
\maketitle
\noindent\textit{Editorial note.} All delivered mathematical text is the author's own, copied from the pinned Stacks Project chapter file homology.tex; no statement or proof step was rewritten, reordered or removed.\par\medskip
\noindent\textit{Editorial note.} The proof is the author's complete argument; both assertions are written out, each with its displayed computation, and the closing sentence is a conclusion rather than a placeholder.\par\medskip
\noindent\textit{Editorial note.} Every labelled reference occurring inside the counted spans is carried in assembly.external\_references as a non-load-bearing navigation pointer (same-chapter ones at their definition site, cross-chapter ones at the invocation inside the counted span), so no load-bearing mathematics is replaced by a stub. No bibliography citation occurs inside the counted spans.\par\medskip
\noindent\textit{Editorial note.} The same-name check for this card scanned every staged Stacks chapter of this repository -- all sources/stacks-*/pinned-original-source.tex files, the card's own chapter excluded -- and the labels found are recorded in dedup.same\_name\_cross\_chapter on the card.\par\medskip
\noindent\textit{Editorial note.} The preamble is the chapter's own native preamble excerpt already used for this source: the comment-type reference, slogan and history environments (native preamble lines 8-13), the Xy-pic setup (native preamble lines 14-18) and the native theorem declarations with the mathematical macros (native preamble lines 154-217, the actual last line of the file).\par\medskip
\noindent\textit{Allowed standard prerequisites (editorial).}\begin{itemize}
\item The connecting morphism of a termwise split sequence. Same chapter, Lemma lemma-ses-termwise-split-cochain, delivered on this card as a context span. Supplies the notation, the hypotheses and the constructed morphism delta whose dependence on the choice of splitting is the subject of the lemma.
\item Homotopy of morphisms of complexes. Same chapter, the definition of a homotopy between two morphisms of complexes. Supplies the meaning of the second assertion, namely the displayed formula expressing delta' minus delta as a boundary.
\end{itemize}\medskip
\begin{lemma}
\label{lemma-ses-termwise-split-cochain}
Let $\mathcal{A}$ be an additive category.
Let
$$
0 \to A^\bullet \to B^\bullet \to C^\bullet \to 0
$$
be a complex (!) of complexes.
Suppose that we are given splittings $B^n = A^n \oplus C^n$
compatible with the maps in the displayed sequence.
Let $s^n : C^n \to B^n$ and $\pi^n : B^n \to A^n$ be the
corresponding maps. Then the family of morphisms
$$
\pi^{n + 1} \circ d_B^n \circ s^n
:
C^n \to A^{n + 1}
$$
define a morphism of complexes $\delta : C^\bullet \to A[1]^\bullet$.
\end{lemma}

\begin{lemma}
\label{lemma-ses-termwise-split-homotopy-cochain}
Notation and assumptions as in
Lemma \ref{lemma-ses-termwise-split-cochain}.
Let $\alpha : A^\bullet \to B^\bullet$,
$\beta : B^\bullet \to C^\bullet$ be the given
morphisms of complexes.
Suppose $(s')^n : C^n \to B^n$ and $(\pi')^n : B^n \to A^n$
is a second choice of splittings.
Write $(s')^n = s^n + \alpha^n \circ h^n$ and
$(\pi')^n = \pi^n + g^n \circ \beta^n$ for some unique
morphisms $h^n : C^n \to A^n$ and $g^n : C^n \to A^n$. Then
\begin{enumerate}
\item $g^n = - h^n$, and
\item the family of maps $\{g^n : C^n \to A[1]^{n - 1}\}$ is a homotopy
between $\delta, \delta' : C^\bullet \to A[1]^\bullet$, more precisely
$(\delta')^n = \delta^n + g^{n + 1} \circ d_C^n + d_{A[1]}^{n - 1} \circ g^n$.
\end{enumerate}
\end{lemma}

\begin{proof}
As $(s')^n$ and $(\pi')^n$ are splittings we have $(\pi')^n \circ (s')^n = 0$.
Hence
$$
0 = ( \pi^n + g^n \circ \beta^n ) \circ ( s^n + \alpha^n \circ h^n ) =
g^n \circ \beta^n \circ s^n + \pi^n \circ \alpha^n \circ h^n =
g^n + h^n
$$
which proves (1). We compute $(\delta')^n$ as follows
$$
( \pi^{n + 1} + g^{n + 1} \circ \beta^{n + 1} )
\circ d_B^n \circ
( s^n + \alpha^n \circ h^n )
= \delta^n + g^{n + 1} \circ d_C^n + d_A^n \circ h^n
$$
Since $h^n = -g^n$ and since $d_{A[1]}^{n - 1} = -d_A^n$ we conclude that (2)
holds.
\end{proof}

\end{document}
